% Section 8.1, from lectures/L08/appendix/a_watchdog.md.

\section{The Watchdog Timer}
\label{c8:sec:watchdog}\label{c8:app:a}

\subsection{The watchdog timer: what it is and why}
\label{c8:a:1}
\begin{itemize}
  \item A watchdog is a second, independent timer with one job: if it isn't reset (``petted'')
    within its configured timeout, it resets the entire chip.
  \item The use cases are:
  \begin{itemize}
    \item A program that might hang.
    \item An infinite loop from a bug.
    \item A sensor that never responds.
    \item A peripheral wait-loop (like \secref{c7:a:2}'s \code{EEPE} wait, or Chapter~\ref{c4:ch}'s
      \code{ADIF} wait) that never sees its flag set.
  \end{itemize}
  \item Without a watchdog, a hang like that freezes the device until someone power-cycles it by
    hand. With one, the chip recovers on its own within one timeout period.
  \item The watchdog timer follows this pattern:
  \begin{itemize}
    \item Enable the watchdog with a timeout comfortably longer than one iteration of
      \code{main()}'s loop normally takes.
    \item Reset the watchdog once per iteration of \code{main()}'s loop.
    \item As long as the loop keeps running, the watchdog never fires. The moment it doesn't
      (because something hung), the watchdog's timeout expires and the chip restarts.
  \end{itemize}
\end{itemize}

\subsection{Watchdog registers}
\label{c8:a:2}
One register, \textbf{\code{WDTCSR}} (Watchdog Timer Control and Status Register), controls both
enablement and timeout via bits \code{WDE} (Watchdog Enable) and \code{WDP0}--\code{WDP3} (Watchdog
Prescaler, selecting the timeout duration).

The four prescaler bits, \code{WDP[3:0]}, set the timeout anywhere from 16\,ms up to 8192\,ms:

\begin{narrowtable}{@{}rcr@{}}
  \thead{Timeout} & \thead{\code{WDP[3:0]}} & \thead{Value} \\
  \midrule
  16\,ms   & \code{0000} & 0  \\
  32\,ms   & \code{0001} & 1  \\
  64\,ms   & \code{0010} & 2  \\
  128\,ms  & \code{0011} & 3  \\
  256\,ms  & \code{0100} & 4  \\
  512\,ms  & \code{0101} & 5  \\
  1024\,ms & \code{0110} & 6  \\
  2048\,ms & \code{0111} & 7  \\
  4096\,ms & \code{1000} & 32 \\
  8192\,ms & \code{1001} & 33 \\
\end{narrowtable}

\noindent The jump from 7 to 32 isn't a typo: \code{WDP2}--\code{WDP0} sit at contiguous bits
2--0 of \code{WDTCSR}, so the first eight timeouts (16\,ms--2048\,ms) map straight onto the
integers 0--7. \code{WDP3}, needed for 4096\,ms and 8192\,ms, lives at bit 5 instead, so those two
timeouts land on \code{1 << 5} (32) and \code{(1 << 5) | (1 << 0)} (33).

This maps cleanly onto a plain \code{enum}: the first eight members can take their default values
(0--7, counting up in declaration order, same as \secref{c1:a:3}'s \code{sensor_state_t}), and
only the last two need an explicit value:

\begin{ccode}
typedef enum
{
    WATCHDOG_TIMEOUT_16MS,        ///< 16 ms timeout.
    WATCHDOG_TIMEOUT_32MS,        ///< 32 ms timeout.
    WATCHDOG_TIMEOUT_64MS,        ///< 64 ms timeout.
    WATCHDOG_TIMEOUT_128MS,       ///< 128 ms timeout.
    WATCHDOG_TIMEOUT_256MS,       ///< 256 ms timeout.
    WATCHDOG_TIMEOUT_512MS,       ///< 512 ms timeout.
    WATCHDOG_TIMEOUT_1024MS,      ///< 1024 ms timeout.
    WATCHDOG_TIMEOUT_2048MS,      ///< 2048 ms timeout.
    WATCHDOG_TIMEOUT_4096MS = 32, ///< 4096 ms timeout.
    WATCHDOG_TIMEOUT_8192MS = 33, ///< 8192 ms timeout.
} watchdog_timeout_t;
\end{ccode}

\noindent Cast a \code{watchdog_timeout_t} directly to \code{uint8_t} and write it to the
\code{WDTCSR} prescaler bits. Its encoded value places the prescaler bits correctly for all ten
timeout settings.

Clearing \code{WDE}, or changing the prescaler bits, needs the same kind of timed sequence as
EEPROM's \code{EEMPE}/\code{EEPE} from Chapter~\ref{c7:ch}, and for the same reason (setting
\code{WDE} alone doesn't, but this book's code always writes \code{WDE} and the timeout together,
so it always uses the sequence):
\begin{itemize}
  \item It's a safety interlock, this time specifically so runaway code can't accidentally disable
    the one thing meant to catch it. \code{WDCE} (Watchdog Change Enable) must be set together with
    \code{WDE}, and the real configuration written within 4 cycles:
\end{itemize}

\begin{ccode}
cli();
WDTCSR |= (1U << WDCE) | (1U << WDE);                      // Unlock the timed sequence.
WDTCSR = (1U << WDE) | (uint8_t)(WATCHDOG_TIMEOUT_1024MS); // Enable watchdog, 1024 ms timeout.
sei();
\end{ccode}

\noindent Resetting the countdown doesn't touch \code{WDTCSR} at all, it's a single dedicated CPU
instruction, \code{WDR} (Watchdog Reset). Nothing about \code{WDR} lives behind a register bit, so
there's no C statement (no \code{WDTCSR |= ...}) that produces it; the only way to emit a specific
machine instruction from C is \textbf{inline assembly}, \code{asm(...)}, which drops a literal
instruction straight into the compiled output, bypassing C entirely for that one line:

\begin{ccode}
asm("WDR");
\end{ccode}

\noindent This is the simplest form GCC supports, ``basic'' inline assembly: just a literal
instruction as a string, no inputs, no outputs, nothing declared as clobbered. (GCC also has
``extended'' inline assembly, with input/output operands and clobber lists, for wiring C variables
into and out of an \code{asm} block; \code{WDR} doesn't read or write anything C can see, so none of
that machinery is needed here.)

One gotcha: \code{asm} is a GNU extension, not standard C, recognized only under a GNU dialect
(\code{-std=gnu11}) or GCC's default mode. This book's demos compile with \code{-std=c11}, strict
ISO C, where the compiler doesn't know the word \code{asm} at all and this line fails to build.
\code{__asm__}, the double-underscore spelling, is always recognized regardless of \code{-std=}, so
that's the form used from here on:

\begin{ccode}
__asm__("wdr");
\end{ccode}

\noindent This book writes straight against the registers instead of reaching for a library, the
same as every peripheral so far (UART, ADC, timers, EEPROM): no header between the code and the
hardware, so nothing about what the chip is actually doing is hidden behind a macro.

\begin{ccode}
#include <avr/interrupt.h>
#include <avr/io.h>

/**
 * @brief Enable the watchdog timer with a 1024 ms timeout.
 */
static void watchdog_enable(void)
{
    cli();
    WDTCSR |= (1U << WDCE) | (1U << WDE);
    WDTCSR = (1U << WDE) | (uint8_t)(WATCHDOG_TIMEOUT_1024MS);
    sei();
}

/**
 * @brief Reset the watchdog timer, restarting its countdown.
 */
static void watchdog_reset(void)
{
    __asm__("wdr");
}
\end{ccode}

\subsection{\code{WDRF}: don't get locked out of disabling the watchdog}
\label{c8:a:3}
A watchdog-caused reset sets \code{WDRF} (Watchdog Reset Flag) in \textbf{\code{MCUSR}} (MCU Status
Register). That flag isn't just informational: the datasheet has \code{WDE} read back as set for as
long as \code{WDRF} is set, no matter what gets written to it. Try to clear \code{WDE} (or disable
the watchdog some other way) while \code{WDRF} is still set, and it silently does nothing, the
watchdog stays armed and fires again next timeout.

This is deliberate, not a bug: it stops a chip that just failed and reset from immediately turning
its own safety net back off, which is exactly the failure mode a watchdog exists to catch. The
practical consequence is that \code{WDRF} must be cleared explicitly, and early, before anything
else touches \code{WDTCSR}:

\begin{ccode}
MCUSR &= ~(1U << WDRF); // Clear the watchdog reset flag; WDE can't be cleared while it's set.
\end{ccode}

\noindent Put this as the first line of \code{setup()}, unconditionally, before
\code{watchdog_enable()} or any other watchdog call. Skip it, and a single watchdog-triggered reset
can leave the chip unable to ever turn the watchdog off again, whatever the running code tries to
do about it.

Clearing \code{WDRF} doesn't stop the watchdog, though, and after a watchdog reset it is still
running. The reset cleared the prescaler bits, so it now runs at its shortest timeout, 16\,ms.
Configure the watchdog straight after clearing the flag, with the timeout you actually want (or turn
it off), before anything slow: a \code{setup()} that first spends 16\,ms on something else, and
printing a single line at 9600 baud already takes that long, is reset again before it gets there,
and again, forever. That's why \secref{c8:a:5}'s \code{setup()} calls \code{watchdog_init()}, which
clears \code{WDRF} and then arms the watchdog, before anything else.

On an Uno you may never see this happen, because its bootloader clears \code{MCUSR} and switches
the watchdog off before it starts your program. Don't rely on that: a program flashed without a
bootloader, or run in a simulator, gets no such help. Nor do Nano boards with the older bootloader
(the ones that need \code{-b 57600}, \secref{c2:d:7}), which after a watchdog reset are known to get
stuck resetting inside the bootloader itself, before your program ever runs.

\subsection{System Reset Mode vs.\ Interrupt Mode}
\label{c8:a:4}
\code{WDE} and \code{WDIE} (Watchdog Interrupt Enable), both in \code{WDTCSR}, pick what a timeout
actually does, independently of each other:
\begin{itemize}
  \item \code{WDE} alone: \textbf{System Reset Mode}. Timeout resets the chip, same as
    \secref{c8:a:2}'s \code{watchdog_enable()}.
  \item \code{WDIE} alone: \textbf{Interrupt Mode}. Timeout jumps to \code{ISR(WDT_vect)} instead
    of resetting anything.
  \item Both set: timeout fires \code{ISR(WDT_vect)} \emph{and} clears \code{WDIE} in hardware on
    the way in, leaving \code{WDE} armed. Pet the watchdog inside that ISR (or soon after) and the
    reset never happens; miss the next timeout too and the chip resets. That's one warning before
    the safety net actually trips.
\end{itemize}
Setting the prescaler bits still needs \secref{c8:a:2}'s \code{WDCE} timed sequence, same as
\code{watchdog_enable()}, just with \code{WDE} left at \code{0} in the value that gets committed.
\code{WDIE} itself isn't behind that sequence, so it's set separately, straight after:

\begin{ccode}
/**
 * @brief Enable the watchdog in Interrupt Mode: ISR(WDT_vect) runs on timeout instead of a
 *        system reset.
 *
 * @param[in] timeout Timeout to arm the watchdog with (see A.2's table).
 */
static void watchdog_enable_interrupt(const watchdog_timeout_t timeout)
{
    watchdog_reset();
    cli();
    WDTCSR |= (1U << WDCE) | (1U << WDE); // Unlock the timed sequence.
    WDTCSR = (uint8_t)(timeout);          // Commit prescaler bits, WDE left at 0.
    sei();
    WDTCSR |= (1U << WDIE);               // Not part of the timed sequence.
}
\end{ccode}

\subsection{Worked example: petting the watchdog in Chapter 7's demo}
\label{c8:a:5}
% Ragged, because the long names in this paragraph leave no good way to justify it.
{\raggedright
Extends \secref{c7:a:4}'s button-toggles-LED1-and-persists-to-EEPROM demo with a 1024\,ms
watchdog (System Reset Mode, \secref{c8:a:2}) on top. \code{setup()} calls
\code{watchdog_init(WATCHDOG_TIMEOUT_1024MS)} once, before anything else, for the reason
\secref{c8:a:3} gives; \code{main()}'s loop calls \code{watchdog_reset()} first thing, every
iteration, whether or not a button event happened that iteration, same as \secref{c8:a:1}'s
pattern. The circuit is the one of Figure~\ref{c6:fig:serial}, unchanged.\par}

\begin{console}
watchdog_driver_demo/
├── Makefile
├── main.c
├── include/
│   └── driver/
│       ├── eeprom.h
│       ├── serial.h
│       └── watchdog.h
└── source/
    └── driver/
        ├── eeprom.c
        ├── serial.c
        └── watchdog.c
\end{console}

\noindent The serial and EEPROM drivers were introduced in Chapters~\ref{c6:ch} and~\ref{c7:ch}
respectively. \file{driver/watchdog.h} and \file{driver/watchdog.c}, printed in full below, are
\secref{c8:a:2}'s \code{watchdog_timeout_t} plus \code{watchdog_init()}/\code{watchdog_reset()},
nothing more.

\subsubsection{The driver in full}
\cfile{lectures/L08/watchdog_driver_demo/include/driver/watchdog.h}
\needspace{8\baselineskip}
\cfile{lectures/L08/watchdog_driver_demo/source/driver/watchdog.c}

\subsubsection{The program}
\code{main.c} in full:

\cfile{lectures/L08/watchdog_driver_demo/main.c}

\noindent Try commenting out just the \code{watchdog_reset();} call and reflashing: nothing about
the rest of the code looks wrong, \code{setup()} still runs, the button still toggles LED1 and
reports over UART, but roughly every 1024\,ms the chip reboots on its own, visible as
\code{setup()}'s LED1-state report firing over and over in the serial monitor even with nobody
touching the button. That's \secref{c8:a:1}'s watchdog doing exactly what it's for: the timeout
wasn't pet in time, so it fired.

This is the one experiment in the book that's simpler to run in the simulator than on the board,
and it's worth doing even if you have hardware: no reflashing between the two versions and no
serial monitor to wire up, so the before-and-after is seconds apart instead of minutes. Petting the
watchdog gives one state report and then silence; starving it gives the same report over and over.
See Appendix~\ref{sim:ch}, which walks through this demo specifically in
\secref{sim:sec:watchdog}.

\subsubsection{Building and running}
Build and flash the same way as \secref{c6:a:6}: on Linux, run \code{make} in the demo's
directory, \file{lectures/L08/watchdog_driver_demo} (\code{make flash} to program the board); on
Windows via Microchip Studio, recreate this folder structure inside the project and add
\file{include} as an include directory. The demo's \file{Makefile} is the one printed in
\secref{c6:sec:build}, with the EEPROM and watchdog drivers added to its source files:

\begin{makecode}
SRC_FILES := main.c \
             source/driver/eeprom.c \
             source/driver/serial.c \
             source/driver/watchdog.c \
\end{makecode}

\subsection{Why no structs here}
\label{c8:a:6}
UART (Chapter~\ref{c6:ch}), EEPROM (Chapter~\ref{c7:ch}) and the watchdog all get away with plain
functions and no struct, unlike where this book is headed next. The ATmega328P has exactly one
UART, exactly one EEPROM and exactly one watchdog; there's no second instance to distinguish
between, no \code{self} parameter needed, nothing to bundle. A struct would just be ceremony around
state that only ever exists once.

LEDs, buttons, and timers aren't so lucky, a program regularly wants several of each, and
Chapter~\ref{c4:ch}'s five-LED exercise (Exercise~\ref{c4:ex:7}) and Chapter~\ref{c5:ch}'s
three-timer exercise (Exercise~\ref{c5:ex:7}) already ran into what that means without a better
tool. That's exactly the gap Chapter~\ref{c9:ch} closes.
